\documentclass{article}
\usepackage{graphicx} % Required for inserting images
\usepackage{tikz}
\usepackage{amsmath, amssymb, amsthm}
\usepackage[most]{tcolorbox}
\usepackage{xcolor}
\usepackage{microtype}

\definecolor{defbg}{RGB}{255, 251, 230}
\definecolor{defborder}{RGB}{230,200,100}
\definecolor{theorembg}{RGB}{230, 251, 255}
\definecolor{theoremborder}{RGB}{100, 200, 230}

\title{Configuration Theory}
\author{Kıvılcım İpek Afet Öztürk}
\date{3 September 2026}

\begin{document}
%% util.tex - utilites for writing papers.

\newcommand{fap}[0]{\,}

\newcommand{\definitionbox}[3]{%
  \begin{tcolorbox}[
      colback=defbg,
      colframe=defborder,
      coltitle=black,
      fonttitle=\bfseries,
      arc=5pt,
      boxrule=0.8pt,
      title={definition #1: #2}
  ]
  \label{#1}
  #3
  \end{tcolorbox}%
}

\newcommand{\theorembox}[3]{%
  \begin{tcolorbox}[
      colback=theorembg,
      colframe=theoremborder,
      coltitle=black,
      fonttitle=\bfseries,
      arc=5pt,
      boxrule=0.8pt,
      title={theorem #1: #2}
  ]
  \label{#1}
  #3
  \end{tcolorbox}%
}

\maketitle

\section{Introduction}
This is a boring paper about how computer systems are configured and how these configurations are shared across clusters of computers.
\\\\
As a very fundamental definition for this paper, we need to define a subset of functions.

\definitionbox{1.1}{\textit{finite function}} {
Let $f: D \to C$ be a function. If $\left|f\right| < \infty$, $f$ is called a \textit{finite function}. 
}

\section{A mathematical model of configuration: dictionaries}
First, we need to take a look at dictionaries, which are the most elementary units of computer configuration.
\definitionbox{2.1}{\textit{dictionary}, \textit{key}, \textit{value}, $\mathcal{K}$, $\mathcal{V}$}{
A \textit{dictionary} is a finite function $f: \mathcal{K} \to \mathcal{V}$ described explicitly by its graph $f = \left\{\left(k_n, v_n\right)\mid \forall{n}\in\mathbb{N}, \left(k_n, v_n\right)\in{\mathcal{K}\times\mathcal{V}}\right\}$.\\\\

Each member of the domain of $f$ is called a \textit{key} of the dictionary.\\
Each member of the codomain of $f$ is called a \textit{value} of the dictionary.\\\\

For the rest of this paper the set of all possible keys a dictionary can have will be denoted as $\mathcal{K}$.\\
For the rest of this paper the set of all possible values a dictionary key may have will be denoted as $\mathcal{V}$.
}

Dictionaries on their own are boring structures, so we need to bring up a few more concepts to make them useful.
\definitionbox{2.2}{\textit{set of software configuration files}, $\mathcal{S}_f$}{
A \textit{set of software configuration files} is a set of files residing in a filesystem that describe the configuration of one instance of software. The parsed set of these configuration files represented as a set of dictionaries, one for each file, is denoted $\mathcal{S}_f$.
}

By knowing these two definitions, we can move on to defining what a configuration means:
\definitionbox{2.3}{\textit{configuration}, $C$}{
    A \textit{configuration} $C$ is the union of a set of software configuration files $S_f$ whose dictionary domains are pairwise disjoint.
    $$C = \bigcup_{f\in S_f} f \,\,\text{where}\,\, \forall{f_i, f_j}\in{S_f} \left(i \neq j \implies \text{Dom}\left(f_i\right)\cap\text{Dom}\left(f_j\right) = \varnothing \right)$$ 
}
\
\section{Describing many pieces of configuration: states}
A computer system may have more than one piece of configuration, because our definitions in \textbf{2.2} and \textbf{2.3} are used for describing configuration for \textit{one} piece of software. A typical computer has more than one pieces of software installed on it. We refer to the whole configuration of a computer as a \textit{state}, which will be our next definition.
\definitionbox{3.1}{state, $\mathcal{S}$}{
    A \textit{state} $\mathcal{S}$ is a set of multiple configurations. States are often denoted as $\mathcal{S}$ and multiple states are denoted with indices below them, such as $\mathcal{S}_1, \mathcal{S}_2, \dots$

    $$\mathcal{S} = \bigcup_{n = 1}^{\left|\mathcal{S}\right|}\left\{C_i\right\}$$
}

For any declarative configuration system; we generally talk about two states: one of them is the desired state of the system, and the other one is the existing state of the system.

\definitionbox{3.2}{desired state, $\mathcal{S_\text{d}}$, $\mathcal{S}$}{
    The intended state of a computer system described by a declarative configuration subsystem is called the \textit{desired state}. Desired state is denoted as $\mathcal{S_\text{d}}$ which is shortened as $\mathcal{S}$ in this paper. 
}
\definitionbox{3.3}{existing state, $\mathcal{S}_\text{ex}$, $\mathcal{E}$}{
    The present state of a computer system described by its configuration files is called the \textit{existing state}. The state is denoted as $\mathcal{S}_\text{ex}$ which is shortened as $\mathcal{E}$ in this paper.
}

\subsection{Methods of comparing states}
Comparing the states is specifically useful when it comes to describing the drift-away of a system from its ideal. States can be compared by a number of different formulas, but this paper will make use of only two of them: set difference between $\mathcal{S}$ and $\mathcal{E}$ and the difference between modification times of the two states.
\definitionbox{3.4}{modification time, $\operatorname{mtime}$}{
    The \textit{modification time} of a configuration file is the last time the configuration file was written to.
    The last write to a configuration file $f$ is denoted as $\mathit{mtime}\,f$. \\\\
    For a configuration $C$, $\operatorname{mtime}C$ is defined as follows:
    $$\operatorname{mtime}C = \max \left\{\operatorname{mtime}f\mid f \subset C\right\}$$
}

\definitionbox{3.5}{symmetric delta, $\operatorname{\Delta}$}{
    For any sets $A$ and $B$, the \textit{symmetric delta} between the two sets is defined as:
    $$A \operatorname{\Delta} B = \left(A\setminus{B}\right)\cup\left(B\setminus{A}\right)$$
}

\definitionbox{3.6}{behind, ahead of, $<$, $>$, drift-away}{
    Let $S_1$ and $S_2$ be two states.\\
    If $\operatorname{mtime}{S_1} < \operatorname{mtime}{S_2}$, $S_1$ is said to be \textit{behind} $S_2$. This is denoted as $S_1 < S_2$ and read as \textit{$S_1$ is behind $S_2$}. \\
    If $\operatorname{mtime}{S_1} > \operatorname{mtime}{S_2}$, $S_1$ is said to be \textit{ahead of} $S_2$. This is denoted as $S_1 > S_2$ and read as \textit{$S_1$ is ahead of $S_2$}.

    The condition of states being behind or ahead of each other is called \textit{drift-away}. 
}
\theorembox{3.7}{existence of a symmetric delta implies drift away}{
    Let $\mathcal{D} = \mathcal{S}\operatorname{\Delta}\mathcal{E}$. If $\mathcal{D} \neq \varnothing$, then $\mathcal{S}$ and $\mathcal{E}$ drift away from each other. 
}
\begin{proof}
    
\end{proof}
\end{document}
